% the N in N-representability problem
% number of modes (number of single particle wavefunctions, local dimension)
% the number of particles
%\def\electron{fermion }
%\def\electrons{fermions }


\documentclass[aps,prl,showpacs,superscriptaddress,floatfix,twocolumn]{revtex4}
\usepackage{amsfonts}
\usepackage{amsmath}

%%%%%%%%%%%%%%%%%%%%%%%%%%%%%%%%%%%%%%%%%%%%%%%%%%%%%%%%%%%%%%%%%%%%%%%%%%%%%%%%%%%%%%%%%%%%%%%%%%%%%%%%%%%%%%%%%%%%%%%%%%%%%%%%%%%%%%%%%%%%%%%%%%%%%%%%%%%%%%%%%%%%%%%%%%%%%%%%%%%%%%%%%%%%%%%%%%%%%%%%%%%%%%%%%%%%%%%%%%%%%%%%%%%%%%
\usepackage{amssymb}
\usepackage{graphics}
\usepackage{graphicx}

%TCIDATA{OutputFilter=LATEX.DLL}
%TCIDATA{Version=5.50.0.2960}
%TCIDATA{<META NAME="SaveForMode" CONTENT="1">}
%TCIDATA{BibliographyScheme=Manual}
%TCIDATA{LastRevised=Friday, November 15, 2024 12:39:43}
%TCIDATA{<META NAME="GraphicsSave" CONTENT="32">}

\newcommand{\be}{\begin{eqnarray}}
\newcommand{\ee}{\end{eqnarray}}
\newcommand{\bea}{\begin{eqnarray}}
\newcommand{\eea}{\end{eqnarray}}
\newcommand{\bma}{\begin{subequations}}
\newcommand{\ema}{\end{subequations}}
\DeclareMathOperator{\poly}{poly}
\DeclareMathOperator{\Tr}{tr}
\newcommand{\bra}[1]{\langle #1 |}
\newcommand{\ket}[1]{| #1 \rangle}
\newcommand{\set}[1]{\lbrace #1 \rbrace}
\newcommand{\RR}{\mathbb{R}}
\newcommand{\SSS}{\mathcal{S}}
\def\N{N}
\def\m{d}
\def\n{N}
\input{tcilatex}

\begin{document}

\title{$N$-representability is QMA-complete}
\author{Yi-Kai \surname{Liu}}
\affiliation{Computer Science and Engineering, University of California, San Diego, US}
\author{Matthias \surname{Christandl}}
\affiliation{Centre for Quantum Computation, Centre for Mathematical Sciences, DAMTP,
University of Cambridge, Wilberforce Road, Cambridge CB3 0WA, United Kingdom}
\author{F. \surname{Verstraete}}
\affiliation{Institute for Quantum Information, Caltech, Pasadena, US}
\affiliation{Facult\"at f\"ur Physik, Universit\"at Wien, Boltzmanngasse 5, A-1090 Wien,
Austria}
\pacs{03.67.�a, 31.25.-v}
\date{\today}

\begin{abstract}
%We show that the $\N$-representability problem, a central object of study in the field of quantum chemistry, is
%QMA-complete and hence NP-hard.
We study the computational complexity of the $N$-representability problem in
quantum chemistry. We show that this problem is QMA-complete, which is the
quantum generalization of NP-complete. Our proof uses a simple mapping from
spin systems to fermionic systems, as well as a convex optimization
technique that reduces the problem of finding ground states to $N$%
-representability. 
%In addition, we show that pure-state $\N$-representability is in QMA,
%and thus is no harder than the standard $\N$-representability problem.
\end{abstract}

\maketitle

The central theoretical problem in the field of many-body strongly
correlated quantum systems is to find efficient ways of simulating
Schr\"odinger's equations: it is very easy to write down those equations,
but notoriously difficult to solve them or even to find approximate
solutions. The main difficulty is the fact that the dimension of the Hilbert
space describing a system of $N$ quantum particles scales exponentially in $N
$. This makes a direct numerical simulation intractable: every time an extra
particle is added to the system, the computational resources would have to
be doubled.

The situation is not hopeless, however, as in principle it could be that all
physical wavefunctions, i.e., the ones that are realized in nature, have
very special properties and can be parameterized in an efficient way. The
idea would then be to propose a variational class of wavefunctions that
capture the physics of the systems of interest, and then do an optimization
over this restricted class. This approach has proven to be very successful,
as witnessed by mean field theory and renormalization group methods. So far,
an efficient variational class to describe complex wavefunctions such as
those arising in quantum chemistry has not been found.

One of the basic problems in quantum chemistry is to find the ground state
of a Hamiltonian describing the many-body system of an atom or molecule. The
essential element that makes typical Hamiltonians very ungeneric is the fact
that at most 2-body interactions occur. This implies that the number of free
parameters in such Hamiltonians scales at most quadratically in the number
of particles or modes, and hence the ground states of all such systems form
a small-dimensional manifold.

In the case of a Hamiltonian with only 2-body interactions, the energy
corresponding to a wavefunction is completely determined by its 2-body
correlation functions, and as a consequence the ground state will be the one
with extremal 2-body reduced density operators. This fact was realized a
long time ago, and led Coulson \cite{Coulson} to propose the following
problem: given a set of $N$ quantum particles, can we characterize the
allowed sets of 2-body correlations or density operators between all pairs
of particles?

If the particles under consideration are fermions, this has been called the 
\textit{$N$-representability problem}~\cite{Coleman}. 
%In the case of fermions, the complete $\N$-body wavefunction can
%always be written as a linear combination of Slater determinants,
%and one is interested in the reduced density operators acting on
%pairs of fermions.
Here, we consider the reduced density operators acting on pairs of fermions,
and we want to decide whether they are consistent with some global state
over $N$ fermions. %A variant of this problem is
%\textit{pure-state} $\N$-representability, where we impose an
%additional requirement that the $\N$ fermions must be in a pure
%state; this problem seems more complicated, because the set of
%feasible reduced density operators is no longer convex.

An efficient solution to the $N$-representability problem would be a huge
breakthrough, as it would (for example) allow us to calculate the binding
energies of all molecules. Therefore, a very large effort has been devoted
to solving this problem and there exists a large literature on this subject
(see e.g. \cite{book1,book,Mazziotti}).

Here we will show that the $N$-representability problem is intractable, as
it is QMA-complete and hence NP-hard. The complexity class QMA (Quantum
Merlin-Arthur) is the natural generalization of the class NP
(nondeterministic polynomial time) to the setting of quantum computing. 
%it was introduced by Kitaev \cite{Kitaev}. %We will also
%show that pure-state $\N$-representability is in QMA, and thus is no
%harder than the standard version of the problem.
Colloquially, a problem is in QMA if there exists an efficient quantum
algorithm that, when given a possible solution to the problem, can verify
whether it is correct; here the ``solution'' may be a quantum state on
polynomially many qubits~\footnote{%
Technically, the problem must be formulated as a yes-or-no question, e.g.,
``for a given Hamiltonian $H$, is the ground state energy less than $\gamma$%
?'' If the answer is yes, there exists a ``witness'' (e.g., the ground
state) which can prove it. Moreover, there is an efficient algorithm that
can check the validity of a witness.}. A problem is QMA-hard if it is at
least as hard as any other problem in QMA; that is, given an efficient
algorithm for this problem, one could solve every other problem in QMA
efficiently~\footnote{%
This includes search and optimization problems, e.g., ``find the ground
state energy of $H$.'' These problems may be QMA-hard, even though they are
not phrased as yes-or-no questions.}. We say that a problem is QMA-complete
if it is in QMA and it is also QMA-hard.

In a seminal work, Kitaev \cite{Kitaev} 
%quantized the proof of Levin and Cook where they showed that 3-SAT is NP-complete and
proved that the Local Hamiltonian problem---determining the ground state
energy of a spin Hamiltonian that is a sum of 5-body terms (on $n$ qubits),
with accuracy $\pm \varepsilon$ where $\varepsilon$ is inverse polynomial in 
$n$---is QMA-complete. In fact, it was later shown that this problem remains
QMA-complete when restricted to 2-body interactions \cite{Kempe1,Kempe2},
and even in the case of geometrically local interactions \cite{Terhal}.

Because the Hamiltonians under consideration are local, the corresponding
ground states have extremal local properties. The dual problem of
determining the ground state energy of a local Hamiltonian is to decide
whether a given set of local density operators can be realized as the
reduced density operators of the same global state. Checking the consistency
of such a set of local density operators $\rho^{(ij)}$, i.e., checking
whether there exists a global state $\sigma$ compatible with those local
density operators, is a QMA-complete problem itself \cite{Yi-Kai}. In the
present paper, we will use very similar techniques to prove that $N$%
-representability, which is the fermionic version of that problem, is also
QMA-complete.

To sketch the main ideas of the proof, we will first consider the classical
marginal problem: %given a probability distribution
%over $\N$ discrete variables that can take the values $\pm 1$, determine the extreme points of all the 2-point
%marginal distributions
%\[p(a_i,a_j)=\sum_{k_1,...k_{i-1},k_{i+1},...k_{j-1},k_{j+1},...}p\left(a_{k_1},a_{k_2},...\right).\]
suppose we have $N$ random variables that can take the values $\pm 1$,
according to some joint probability distribution, and we define the
2-variable marginal distributions 
\[
p(a_i,a_j) = \sum_{a_1,...,a_{i-1},a_{i+1},...,a_{j-1},a_{j+1},...,a_N}
p(a_1,a_2,...,a_N). 
\]
Given a set of marginal probability distributions, we would like to check
whether they are consistent, i.e., whether there exists a global probability
distribution whose marginals are equal to the ones we were given.

Suppose that there exists an efficient algorithm (whose running time is
polynomial in $N$) to solve this problem. Then it would be possible to
identify ground states of all Ising spin glasses. The strategy would be as
follows: since the set of consistent marginals is convex and the energy is a
linear function of the marginals, the problem amounts to minimizing a linear
function subject to a set of convex constraints that can be checked in
polynomial time. This can be done in polynomial time using the ellipsoid
method \cite{GLS}. It therefore follows that an efficient algorithm for
checking consistency allows to find ground states of Ising spin glasses, a
problem that is known to be NP-hard \cite{Barahona}. Hence the problem of
determining whether a set of binary marginals is consistent is itself
NP-hard. 
% Note: finding ground states of Ising models is not in NP because it's a search problem, not a decision problem.
% Also, consistency of marginals is not in NP because it's not clear how to represent the global probability
% distribution as a classical witness; however, the problem is in MA.

Let us now formulate the $N$-representability problem in the context of
quantum chemistry. Electrons and nuclei tend to arrange themselves such as
to minimize their energy, and the binding energy of a molecule can be
determined by calculating the minimal energy of the corresponding
Hamiltonian. In practice, the nuclei are fairly well localized and can be
treated as classical degrees of freedom, and the wavefunction of the $N$
electrons can be approximated as a linear combination of tensor products of
the $d$ single-particle modes in the system (those form the basis set).

As electrons are fermionic, the complete wavefunction must be antisymmetric,
and this is most easily taken into account by working in the formalism of
second quantization: 
\[
|\psi\rangle = \sum_{{\tiny 
\begin{array}{l}
i_1,...,i_d=0 \\ 
i_1+...+i_d=N%
\end{array}%
}}^1 c_{i_1,...,i_d} (a_1^\dagger)^{i_1}...(a_d^\dagger)^{i_d}
|\Omega\rangle. 
\]
Here $a_j^\dagger$ is the creation operator for the $j$'th mode, and $%
|\Omega\rangle$ represents the vacuum state without fermions. The creation
and annihilation operators obey the following anticommutation relations: 
\[
\{a_i,a_j\} = 0 = \{a_i^\dagger,a_j^\dagger\} \hspace{1cm}
\{a_i,a_j^\dagger\} = 2\delta_{ij}. 
\]
Note that we restrict ourselves to the subspace of states with exactly $N$
fermions. $d$ denotes the number of modes, which is typically much larger
than $N$. The number of degrees of freedom is $\binom{d}{N}$, which grows
exponentially in $N$ when $d \geq cN$ for some constant $c>1$.

In the case of quantum chemistry, the Hamiltonian typically contains only
one- and two-body interactions between all modes, and so it can be written
as a linear combination of terms of the form $a_i^\dagger a_j$ and $%
a_i^\dagger a_j^\dagger a_l a_k$. 
%and the energy is therefore a linear combination of terms
%\bea \rho^{(1)}_{ij}&\equiv&\langle\psi|a_i^\dagger a_j|\psi\rangle\\
%\rho^{(2)}_{ijkl}&\equiv&\langle\psi|a_i^\dagger a_j^\dagger a_ka_l|\psi\rangle\label{cons}\eea

The 2-fermion reduced density matrix (2-RDM) is calculated by tracing out
all but two of the fermions: 
\[
\rho^{(2)} = \Tr_{3,\ldots,N} \rho^{(N)}. 
\]
where $\rho^{(N)}$ is a mixture of states $| \psi \rangle$ with exactly $N$
fermions. In the language of second quantization, the matrix elements of the
2-RDM are given by: 
\[
\rho^{(2)}_{ijkl} = \frac{1}{N(N-1)} \langle a_k^\dagger a_l^\dagger a_j a_i
\rangle. 
\]

The $N$-representability problem (with $d$ modes) can now be stated as
follows. Consider a system of $N$ fermions and $d$ modes, $d \leq \poly(N)$.
(For purposes of complexity, we consider $N$ to be the ``size'' of the
problem.) We are given a 2-fermion density matrix $\rho$, of size $\frac{%
d(d-1)}{2} \times \frac{d(d-1)}{2}$, where each entry is specified with $%
\poly(N)$ bits of precision. In addition, we are given a real number $\beta
\geq 1/\poly(N)$, specified with $\poly(N)$ bits of precision. The problem
is to distinguish between the following two cases:

\begin{itemize}
\item There exists an $N$-fermion state $\sigma$ such that $\Tr%
_{3,\ldots,N}(\sigma) = \rho$. In this case, answer ``YES.''

\item For all $N$-fermion states $\sigma$, $\lVert \Tr_{3,\ldots,N}(\sigma)
- \rho \rVert_1 \geq \beta$. In this case, answer ``NO.''
\end{itemize}

If neither of these cases applies, then one may answer either ``YES'' or
``NO.'' Note that we do not insist on solving the problem exactly; we allow
an error tolerance of $\beta \geq 1/\poly(N)$. We use the $\ell_1$ matrix
norm or trace distance, $\lVert A \rVert_1 = \Tr |A|$, to measure the
distance between $\sigma$ and $\rho$.

%We define the pure-state $\N$-representability problem in a similar
%way, with the following modifications.  In addition to $\rho$ and
%$\beta$, we are given a real number $\delta \geq 1/\poly(\n)$,
%specified with $\poly(\n)$ bits of precision.  We have to distinguish
%between these two cases:
%\begin{itemize}
%\item There exists an $\n$-fermion pure state $\sigma$ such that
%$\Tr_{3,\ldots,\n}(\sigma) = \rho$.  In this case, answer ``YES.''
%\item For all $\n$-fermion states $\sigma$ such that $\Tr(\sigma^2)
%\geq 1-\delta$, we have that
%$\lVert \Tr_{3,\ldots,\n}(\sigma) - \rho \rVert_1 \geq \beta$.
%In this case, answer ``NO.''
%\end{itemize}
%Note that we use $\Tr(\sigma^2)$ to measure the purity of the state
%$\sigma$, and we allow an error tolerance $\delta \geq 1/\poly(\n)$.

We will show that $N$-representability is QMA-complete. The proof consists
of two parts. First, we show that any 2-local Hamiltonian of spins can be
simulated using a 2-local Hamiltonian of fermions with $d=2N$. Using
techniques of convex programming, we show that an oracle for $N$%
-representability would allow us to estimate the ground state energies of
2-local Hamiltonians; thus, $N$-representability is QMA-hard.

Second, we show that $N$-representability is in QMA; specifically, we
construct a quantum verifier that can check whether a 2-particle state is $N$%
-representable, given a suitable witness. 
%by using the fact that QMA+=QMA \cite{Aharonov} and that consistency of
%local density operators of a fermionic state of $\m$ modes can be checked by calculating many-body expectation values
%on a $N$-qubit state obtained by performing a Jordan-Wigner transformation.
%We also show that pure-state $\N$-representability is in QMA, using similar techniques.

\vskip 10pt

Let us first show how to map a 2-local Hamiltonian, defined on a system of $N
$ qubits, to a 2-local Hamiltonian on fermions, with $d = 2N$ modes; this is
the opposite of what has been done in \cite{VC}. The idea is to represent
each qubit $i$ as a single fermion that can be in two different modes $%
a_i,b_i$; so each $N$-qubit basis state corresponds to the following $N$%
-fermion state: 
\begin{equation}
| z_1 \rangle \otimes \cdots \otimes | z_N \rangle \mapsto
(a_1^\dagger)^{1-z_1} (b_1^\dagger)^{z_1} \cdots (a_N^\dagger)^{1-z_N}
(b_N^\dagger)^{z_N} | \Omega \rangle.  \label{map0}
\end{equation}
Also, all the relevant single-qubit operators should correspond to bilinear
functions of the creation and annihilation operators (this construction
guarantees that operators on different qubits commute). This can be achieved
by the following mapping: 
\begin{equation}
\sigma^x_i \leftrightarrow a_i^\dagger b_i + b_i^\dagger a_i, \ \sigma^y_i
\leftrightarrow i\left( b_i^\dagger a_i - a_i^\dagger b_i \right), \
\sigma^z_i \leftrightarrow \openone - 2 b_i^\dagger b_i.  \label{map}
\end{equation}
%\begin{eqnarray}
%\sigma^x_i &\leftrightarrow& a_i^\dagger b_i + b_i^\dagger a_i \nonumber\\
%\sigma^y_i &\leftrightarrow& i\left( b_i^\dagger a_i - a_i^\dagger b_i \right) \nonumber\\
%\sigma^z_i &\leftrightarrow& \openone - 2 b_i^\dagger b_i \label{map}
%\end{eqnarray}
It can easily be checked that these fermionic operators obey the Pauli
commutation relations within the subspace spanned by states of the form (\ref%
{map0}).

Recall that the qubit Hamiltonian can be written as a linear combination of
Pauli operators. So the fermionic Hamiltonian can now be created by
rewriting these Pauli operators with respect to the above mapping (note that
only bilinear terms occur, and hence operators on different sites commute).
Notice that if the qubit Hamiltonian only contains 2-body terms, then the
fermionic Hamiltonian will only contain terms with at most 2 annihilation
and 2 creation operators.

The only thing that is left to do is to guarantee that there is exactly one
fermion on every site $i$. This can be achieved by adding the following
projectors as extra terms in the fermionic Hamiltonian: 
\[
P_i=(2a_i^\dagger a_i-\openone)(2b_i^\dagger b_i-\openone).
\]
All the $P_i$ are biquadratic and commute with all the operators introduced
in (\ref{map}), and hence the complete Hamiltonian will be block diagonal.
By making the weights of these projectors large enough (a constant times the
norm of the Hamiltonian, which is at most polynomial in $N$), we can always
guarantee that the ground state of the full Hamiltonian will have exactly
one fermion per site \footnote{%
Another option would be to represent each qubit $i$ using either zero or two
fermions, e.g., for a single qubit, $| 0 \rangle$ and $| 1 \rangle$
correspond to $| \Omega \rangle$ and $a_i^\dagger b_i^\dagger | \Omega
\rangle$. Define the mapping $\sigma^x_i \leftrightarrow \left( a_i^\dagger
b_i + b_i^\dagger a_i \right) +  \left( b_i a_i + a_i^\dagger b_i^\dagger
\right)$, $\sigma^y_i \leftrightarrow i\left( b_i^\dagger a_i - a_i^\dagger
b_i \right) +  i\left( a_i^\dagger b_i^\dagger - b_i a_i \right)$, $%
\sigma^z_i \leftrightarrow \openone - 2 b_i^\dagger b_i$. Then the
Hamiltonian would act identically on the subspace with one fermion per site,
and on the subspace with zero or two fermions per site, 
%Then the subspace with even fermion occupation number would be isomorphic to the one with odd occupation number,
and even a small constant in front of the projectors $P_i$ would be enough
to guarantee that we end up in the right subspace.}.

\vskip 10pt

Let us now assume that we have an efficient algorithm for $N$%
-representability. We claim that this allows us to efficiently determine the
ground state energy of a 2-local Hamiltonian on qubits, a problem which is
known to be QMA-hard \cite{Kempe2}. We start by transforming the Hamiltonian 
$H_{\text{qubit}}$ on $N$ qubits with 2-particle interactions into a
Hamiltonian $H_{\text{fermi}}$ on $d=2N$ fermionic modes with 2-particle
interactions (as described above). The problem of finding the ground state
energy of $H_{\text{fermi}}$ can be expressed as a convex program, which can
be solved in polynomial time using an oracle for $N$-representability (this
is what we will show below). Our approach is similar to \cite{Yi-Kai},
though in this case we are dealing with fermions rather than qubits.

The basic idea is to construct a convex program that finds a 2-particle
density matrix $\rho$ which is $N$-representable, and which minimizes the
expectation value of $H_{\text{fermi}}$. This program has polynomially many
variables, and it is easy to see that the set of $N$-representable states is
convex, and $\langle H_{\text{fermi}} \rangle$ is a linear function of $\rho$%
.

There are two minor complications. First, $H_{\text{fermi}}$ describes a
system with $2N$ modes and an arbitrary number of particles, whereas $N$%
-representability pertains to a system with exactly $N$ particles. But this
is not a problem, because the interesting behavior in $H_{\text{fermi}}$
occurs in a subspace where all the states have exactly $N$ particles. So, in
our convex program, we only need to consider states $\rho$ that have exactly 
$N$ particles.

This lets us simplify $H_{\text{fermi}}$ as follows. Since we are only
interested in how it acts on $N$-particle states, we can write $a_i^\dagger
a_j = \frac{1}{N-1} a_i^\dagger (\sum_k a_k^\dagger a_k) a_j$. So we can
assume that all the terms in $H_{\text{fermi}}$ are of the form $a_i^\dagger
a_j^\dagger a_l a_k$.

Second, convex optimization algorithms usually require that the set $K$ of
feasible solutions be full-dimensional, i.e., $K$ cannot lie in a
lower-dimensional subspace. So we have to represent the state $\rho$ in such
a way that there are no redundant variables.

Let $\mathcal{S}$ be a complete set of 2-particle observables \footnote{%
One possible set of observables is the following: First, define $a_I =
a_{i_2} a_{i_1}$, for all pairs $I = \lbrace i_1,i_2 \rbrace$, $i_1<i_2$.
Also fix an ordering on the pairs $I$. We now define the following
observables: $X_{IJ} = a_I^\dagger a_J + a_J^\dagger a_I$, for all $I \prec J
$; $Y_{IJ} = -i a_I^\dagger a_J +i a_J^\dagger a_I$, for all $I \prec J$;
and $Z_I = a_I^\dagger a_I$, for all $I$ except the last one. These
operators are Hermitian, with eigenvalues in the interval $[-1,1]$. Taking
real linear combinations, these operators form a basis for the space of
2-fermion density matrices.}, and let $\ell = |\mathcal{S}|$; note that $%
\ell \leq \poly(d) \leq \poly(N)$. We represent $\rho$ in terms of its
expectation values $\alpha_S = \Tr(S\rho)$ for all $S \in \mathcal{S}$; let $%
\vec{\alpha} \in \mathbb{R}^\ell$ denote the vector of these expectation
values. We define $K$ to be the set of all $\vec{\alpha}$ such that the
corresponding state $\rho$ is $N$-representable. Note that the $N$%
-representability oracle allows us to test whether a given point $\vec{\alpha%
}$ is in $K$.

We write our Hamiltonian in the form $H = \sum_{S \in \mathcal{S}} \gamma_S S
$ (plus a constant term); the coefficients $\gamma_S$ can be computed using
the Gram-Schmidt procedure. It is easy to see that $\Tr(H\rho) = \sum_{S \in 
\mathcal{S}} \gamma_S \alpha_S$. Our convex program is as follows: find some 
$\vec{\alpha} \in K$ that minimizes $f(\vec{\alpha}) = \sum_{S \in \mathcal{S%
}} \gamma_S \alpha_S$.

We can solve this convex program in polynomial time, using the shallow-cut
ellipsoid algorithm (see \cite{GLS}, and references cited therein). We
mention a few technical details. The algorithm requires some additional
information about $K$, namely a guarantee that $K$ is contained in a ball of
radius $R$ centered at 0, and $K$ contains a ball of radius $r$ centered at
some point $p$. (Also, the running time of the algorithm grows polynomially
in $\log(R/r)$.) In our case, we can set $R = \sqrt{\ell}$, and $r = 1/\poly%
(\ell)$ \footnote{%
It is easy to see that $K$ is contained in a ball of radius $R = \sqrt{\ell}$%
, since for all $\vec{\alpha} \in K$, we have $-1 \leq \alpha_S \leq 1$. 
%____old version____
%\vskip 1pt \hspace{4pt}
%The second claim, that $K$ contains a
%ball of radius $r = 1/\poly(\ell)$, is less trivial.  We need to show that,
%starting from 0 (i.e., the fully mixed state on 2 fermions), if we add a
%$1/\poly(\ell)$ perturbation in any direction, the resulting state will
%still be $\N$-representable.
%\vskip 1pt \hspace{4pt}
%First consider the classical version of the problem, where we only
%consider the diagonal observables $Z_I$, and we restrict ourselves to
%classical mixtures of basis states.  Fix the number of modes $\m$.
%The claim is true when we have $\n=2$ particles, and likewise when we
%have $\n=\m-2$ particles.  Moreover, when we decrease the number of
%particles from $\n$ to $\n-1$ (for instance by choosing one of the
%particles at random and throwing it away), the ball contained in $K$
%shrinks by a factor of $1-\frac{2}{\n}$ at worst.  Thus the claim
%remains true when we have $\n=\m/2$ particles.
%\vskip 1pt \hspace{4pt}
%Now we return to the original quantum problem.  Take an arbitrary
%$1/\poly(\ell)$ perturbation away from 0.  We will show that there is
%an $\n$-particle state with these expectation values.  From the above
%argument, we know there is an $\n$-particle diagonal state $\sigma_{\text{dg}}$
%that has the desired $Z_I$ expectation values.  We now perturb
%$\sigma_{\text{dg}}$ by adding $X_{IJ}$ and $Y_{IJ}$ terms, to get
%the desired expectation values for these observables.  Note that the
%$X_{IJ}$ and $Y_{IJ}$ are orthogonal, so these perturbations do not
%interfere with each other or with $\sigma_{\text{dg}}$.  Also,
%without loss of generality, we can assume that $\sigma_{\text{dg}}$
%has all its eigenvalues bounded away from 0, say by $(2\binom{\m}{\n})^{-1}$,
%so that after the perturbations, the state is still positive semidefinite.
%____old version____
%\vskip 1pt \hspace{4pt}
%The second claim, that $K$ contains a
%ball of radius $r = 1/\poly(\ell)$, is less trivial.  The proof is as
%follows.  We will consider $\N$-representability for different values
%of $\N$; let $K_{\n}$ denote the set of all vectors $\vec{\alpha}$ that
%are $\n$-representable.  Obviously, $K_2$ contains a ball of radius
%$1/\poly(\ell)$ (this is the trivial case).  We will show that $K_{\m-2}$
%also contains a ball of radius $1/\poly(\ell)$, where ``$\poly$'' here
%refers to a different polynomial (proof deferred to the next paragraph).
%Also, notice that $K_{\n-1} \supseteq K_{\n}$, for all $3 \leq \n \leq
%\m-2$ (since, if $\vec{\alpha}$ is consistent with some $\n$-fermion
%state $\sigma$, then it is consistent with the $(\n-1)$-fermion state
%$\sigma'$ that results from tracing out the last particle).  The claim
%then follows immediately.
%\vskip 1pt \hspace{4pt}
%Now we will prove that $K_{\m-2}$ contains a ball of radius
%$1/\poly(\ell)$.  First, consider what happens when we switch the roles
%of the annihilation and creation operators, i.e., we replace $a_i$ with
%$a_i^\dagger$, and vice versa.  Then $a_I^\dagger a_J$ becomes $a_I
%a_J^\dagger$.  These are the ``particle-hole'' operators.  We restrict
%them to act on the subspace of $\n$-fermion states; then we can write the
%``particle-hole'' operators as linear combinations of the standard
%operators, and vice versa; e.g., if $I = \set{i_1,i_2}$, then $a_I
%a_I^\dagger = 1 - a_{i_1}^\dagger a_{i_1} - a_{i_2}^\dagger a_{i_2} +
%a_I^\dagger a_I = \binom{\n}{2}^{-1} \sum_J a_J^\dagger a_J - (\n-1)^{-1}
%\sum_{i_1 \in J} a_J^\dagger a_J - (\n-1)^{-1} \sum_{i_2 \in J}
%a_J^\dagger a_J + a_I^\dagger a_I$; similar but more complicated
%equations hold in the other cases.
%\vskip 1pt \hspace{4pt}
%Motivated by the above, we define ``particle-hole'' observables:
%$X'_{IJ} = a_I a_J^\dagger + a_J a_I^\dagger$; $Y'_{IJ} = -i a_I
%a_J^\dagger +i a_J a_I^\dagger$; and $Z'_I = a_I a_I^\dagger$.  Let
%$\vec{\alpha}'$ denote a vector containing expectation values for these
%observables; let $K'_{\n}$ be the set of all $\vec{\alpha}'$ that are
%$\n$-representable.  We know that we can write the ``particle-hole''
%observables as linear combinations of the standard observables, and vice
%versa.  Thus there is a linear transformation $A$ that maps $K_{\n}$ to
%$K'_{\n}$, and moreover, $A$ is invertible.
%\vskip 1pt \hspace{4pt}
%In particular, $A$ maps $K_2$ to $K'_2$; and, because of the symmetry
%between particles and ``particle holes,'' we have that $K'_2 = K_{\m-2}$.
%Thus $A$ maps $K_2$ to $K_{\m-2}$.  Now we are almost done; we know that
%$K_2$ contains a ball of radius $1/\poly(\ell)$, and we just need to show
%that the map $A$ does not shrink the radius of the ball too much.  Write
%the singular value decomposition $A = UDV$, where $U$ and $V$ are
%unitary, and $D$ is diagonal, with diagonal entries $D_{ii}>0$.  Let
%$B = A^{-1}$.  We can calculate the matrix elements of $B$, and show
%that $\Tr(B^\dagger B) = \sum_{ij}
%|B_{ij}|^2 \leq \poly(\ell)$.  At the same time, $\Tr(B^\dagger B) =
%\Tr(D^{-2}) \geq D_{ii}^{-2}$, for all $i$.  So we have $D_{ii} \geq
%1/\poly(\ell)$, for all $i$.  This implies that the map $A$ shrinks the
%radius of the ball by at most a $\poly(\ell)$ factor.  So $K_{\m-2}$
%contains a ball of radius $1/\poly(\ell)$.
%____current version____
\vskip 1pt \hspace{4pt} The second claim, that $K$ contains a ball of radius 
$r = 1/\poly(\ell)$, is less trivial. The proof is as follows. \vskip 1pt 
\hspace{4pt} We will consider $N$-representability for different values of $N
$; let $K_{N}$ denote the set of all vectors $\vec{\alpha}$ that are $N$%
-representable. Obviously, $K_2$ contains a ball of radius $1/\poly(\ell)$
(this is the trivial case). We will prove that $K_{N}$ contains a ball of
radius $1/\poly(\ell)$, for all $2 \leq N \leq d-2$. \vskip 1pt \hspace{4pt}
We define ``particle-hole'' observables, by replacing $a_i$ with $a_i^\dagger
$, and vice versa: $X^{\prime }_{IJ} = a_I a_J^\dagger + a_J a_I^\dagger$; $%
Y^{\prime }_{IJ} = -i a_I a_J^\dagger +i a_J a_I^\dagger$; and $Z^{\prime
}_I = a_I a_I^\dagger$. Let $\vec{\alpha}^{\prime }$ denote a vector
containing expectation values for these observables; let $K^{\prime }_{N}$
be the set of all $\vec{\alpha}^{\prime }$ that are $N$-representable. %
\vskip 1pt \hspace{4pt} First, we claim that $K_2 = K^{\prime }_{d-2}$.
Notice that there is a natural correspondence between the 2-particle Slater
basis states and the $(d-2)$-particle Slater basis states: the 2-particle
state with modes $i$ and $j$ occupied corresponds to the $(d-2)$-particle
state with modes $i$ and $j$ empty. \vskip 1pt \hspace{4pt} So take any
point $\alpha \in K_2$, which represents the expectation values of the
2-particle observables for some 2-particle state $\sigma$. Use $\sigma$ to
construct a $(d-2)$-particle state $\tau$, by replacing each 2-particle
Slater basis state with the corresponding $(d-2)$-particle Slater basis
state. Then the expectation values of the 2-particle observables for $\sigma$
are exactly the expectation values of the 2-hole observables for $\tau$. So $%
\alpha$ is in $K^{\prime }_{d-2}$. This shows that $K_2 \subseteq K^{\prime
}_{d-2}$. A similar argument shows that $K^{\prime }_{d-2} \subseteq K_2$.
So $K_2 = K^{\prime }_{d-2}$. \vskip 1pt \hspace{4pt} Next, we claim that
there is an invertible linear transformation $A$ that maps $K^{\prime }_{d-2}
$ to $K_{d-2}$. \vskip 1pt \hspace{4pt} First we map $K^{\prime }_{d-2}$ to $%
K_{d-2}$. Observe that we can write each 2-particle operator as a linear
combination of 2-hole operators. (This holds provided we restrict the
operators to act only on the subspace of $(d-2)$-particle states.) For
instance, when $I \cap J = \emptyset$, $a_I^\dagger a_J = a_J a_I^\dagger$.
When $|I \cap J| = 1$, we write equations such as $a_i^\dagger a_l^\dagger
a_l a_j  = a_i^\dagger a_j a_l^\dagger a_l  = -a_j a_i^\dagger (1 - a_l
a_l^\dagger)  = -a_j a_i^\dagger + a_j a_i^\dagger a_l a_l^\dagger  = -a_j
a_i^\dagger + a_l a_j a_i^\dagger a_l^\dagger$; then use the fact that $a_j
a_i^\dagger  = (\sum_{k \neq i,j} a_k a_k^\dagger) a_j a_i^\dagger  =
\sum_{k \neq i,j} a_k a_j a_i^\dagger a_k^\dagger$. When $I = J$, we have $%
a_I^\dagger a_I = \sum_{L \cap I = \emptyset} a_L a_L^\dagger$. (For each of
these equations, it is easy to check that the left and right sides act
identically on all $(d-2)$-particle Slater basis states.) \vskip 1pt \hspace{%
4pt} Thus, for any $(d-2)$-particle state $\sigma$, the expectation values
of the 2-particle observables are linear functions of the expectation values
of the 2-hole observables. Thus we have a linear transformation $A$ that
maps $K^{\prime }_{d-2}$ to $K_{d-2}$. \vskip 1pt \hspace{4pt} Similarly, we
can map $K_{d-2}$ to $K^{\prime }_{d-2}$. We write each 2-hole operator as a
linear combination of 2-particle operators (again, restricting the operators
to act only on the subspace of $(d-2)$-particle states). For instance, when $%
I \cap J = \emptyset$, $a_I a_J^\dagger = a_J^\dagger a_I$. When $|I \cap J|
= 1$, we write equations such as $a_l a_i a_j^\dagger a_l^\dagger  = a_i
a_j^\dagger a_l a_l^\dagger  = -a_j^\dagger a_i (1 - a_l^\dagger a_l)  =
-a_j^\dagger a_i + a_j^\dagger a_i a_l^\dagger a_l  = -a_j^\dagger a_i +
a_j^\dagger a_l^\dagger a_l a_i$; then use the fact that $a_j^\dagger a_i  =
(\frac{1}{d-3} \sum_{k \neq i,j} a_k^\dagger a_k) a_j^\dagger a_i  = \frac{1%
}{d-3} \sum_{k \neq i,j} a_j^\dagger a_k^\dagger a_k a_i$. When $I = J$, we
write $a_I a_I^\dagger  = a_{i_1} a_{i_1}^\dagger a_{i_2} a_{i_2}^\dagger  =
(1 - a_{i_1}^\dagger a_{i_1}) (1 - a_{i_2}^\dagger a_{i_2})  = 1 -
a_{i_1}^\dagger a_{i_1} - a_{i_2}^\dagger a_{i_2} + a_I^\dagger a_I$; then
use the fact that $1 = \binom{d-2}{2}^{-1} \sum_L a_L^\dagger a_L$, and $%
a_i^\dagger a_i  = \frac{1}{d-3} \sum_{l \neq i} a_l^\dagger a_l a_i^\dagger
a_i  = \frac{1}{d-3} \sum_{l \neq i} a_i^\dagger a_l^\dagger a_l a_i$. %
\vskip 1pt \hspace{4pt} Thus, for any $(d-2)$-particle state $\sigma$, the
expectation values of the 2-hole observables are linear functions of the
expectation values of the 2-particle observables. Thus we have a linear
transformation $B$ that maps $K_{d-2}$ to $K^{\prime }_{d-2}$, and $B =
A^{-1}$. \vskip 1pt \hspace{4pt} We claim that $K_{d-2}$ contains a ball of
radius $1/\poly(\ell)$. We know that $K^{\prime }_{d-2}$ contains a ball of
radius $1/\poly(\ell)$ (since $K_2 = K^{\prime }_{d-2}$), and we will show
that the map $A$ does not shrink this too much. Write the singular value
decomposition $A = UDV$, where $U$ and $V$ are unitary, and $D$ is diagonal,
with diagonal entries $D_{ii}>0$. Let $B = A^{-1}$. Looking at the matrix
elements of $B$, we can see that $\Tr(B^\dagger B) = \sum_{ij} |B_{ij}|^2
\leq \poly(\ell)$. At the same time, $\Tr(B^\dagger B) = \Tr(U D^{-1} V
V^{-1} D^{-1} U^{-1}) = \Tr(D^{-2}) \geq D_{ii}^{-2}$, for all $i$. So we
have $D_{ii} \geq 1/\poly(\ell)$, for all $i$. Thus the map $A$ shrinks the
radius of the ball by at most a $\poly(\ell)$ factor. So $K_{d-2}$ contains
a ball of radius $1/\poly(\ell)$. \vskip 1pt \hspace{4pt} Now we are almost
done. Notice that $K_{N-1} \supseteq K_{N}$, for all $3 \leq N \leq d-2$
(since, if $\vec{\alpha}$ is consistent with some $N$-fermion state $\sigma$%
, then it is consistent with the $(N-1)$-fermion state $\sigma^{\prime }$
that results from tracing out the last particle). So we conclude that $K_{N}$
contains a ball of radius $1/\poly(\ell)$, for all $2 \leq N \leq d-2$. This
completes the proof.}.

There is also the issue of numerical precision. Our oracle for $N$%
-representability has limited accuracy---it may give incorrect answers for
points near the boundary of $K$. On the other hand, we only need an
approximate solution to the convex program, in order to solve the Local
Hamiltonian problem. It turns out that $1/\poly(N)$ precision is sufficient.
The ellipsoid algorithm still works in this setting; see \cite{GLS} for
details.

Note that, in place of the ellipsoid algorithm, we could have used a
different algorithm based on random walks in convex bodies \cite{BV}; this
was the approach used in \cite{Yi-Kai}. However, it is not clear if we could
use one of the interior-point methods for convex optimization; these methods
are substantially faster, but they usually require an explicit description
of the constraints, not just a membership oracle.

This completes the proof that $N$-representability is QMA-hard. As a
corollary, we have also proven that estimating the ground state energy for
local fermionic Hamiltonians is QMA-hard. Note that the use of the ellipsoid
algorithm to reduce a convex optimization problem to a convex membership
problem is not new; see \cite{GLS} for references to related work.

\vskip 10pt

Next, we show that $N$-representability is in QMA. That is, we construct a
poly-time quantum verifier $V$ that takes two inputs: a description of the
problem (that is, $\rho$ and $\beta$); and a ``witness'' $\tau$, which is a
quantum state on polynomially many qubits. The verifier $V$ should have the
following property: if $\rho$ is $N$-representable, there exists a witness $%
\tau$ that causes $V$ to output ``true'' with probability $\geq p_1$; if $%
\rho$ is not $N$-representable (within error tolerance $\beta$), then for
all possible states $\tau$, $V$ outputs ``true'' with probability $\leq p_0$%
; and $p_1 - p_0 \geq 1/\poly(N)$.

The idea is that, when $\rho$ is $N$-representable, the correct witness $\tau
$ consists of (multiple copies of) an $N$-fermion state $\sigma$ that
satisfies $\Tr_{3,\ldots,N}(\sigma) = \rho$. Then the verifier can use
quantum state tomography to compare $\sigma$ and $\rho$.

%Say we have a 2-fermion state $\rho$ which is $\N$-representable;
%then it is natural to use the corresponding $\n$-fermion state
%$\sigma$ as the witness.

We represent the $N$-fermion state $\sigma$ using $d$ qubits, via the
following mapping: 
\[
(a_1^\dagger)^{i_1} ... (a_d^\dagger)^{i_d} | \Omega \rangle \leftrightarrow
| i_1 \rangle \otimes ... \otimes | i_d \rangle. 
\]
Call the resulting qubit state $\tilde{\sigma}$. We use the Jordan-Wigner
transform to map the fermionic annihilation operators to qubit operators: 
\[
a_i \leftrightarrow A_i \equiv -(\otimes_{k<i}\hspace{.1cm}\sigma_k^z)
\otimes | 0 \rangle \langle 1 |_i. 
\]
Thus, an observable $O = a_i^\dagger a_j^\dagger a_l a_k + a_k^\dagger
a_l^\dagger a_j a_i$ is transformed into $\tilde{O} = A_i^\dagger
A_j^\dagger A_l A_k + A_k^\dagger A_l^\dagger A_j A_i$, which is a tensor
product of many single-qubit observables and one four-qubit observable.

We claim that the expectation value $\langle \tilde{O} \rangle$ can be
estimated efficiently. Without loss of generality, assume that the
eigenvalues of $\tilde{O}$ lie in the interval $[0,1]$. Then there is an
efficiently-implementable measurement which outputs ``1'' with probability $%
\langle \tilde{O} \rangle$ \footnote{%
For example, let $\lambda_i$ and $| \theta_i \rangle$ be the eigenvalues and
eigenvectors of $\tilde{O}$, and note that they are easy to compute. Add an
ancilla qubit, perform the unitary operation $U:\: | \theta_i \rangle| 0
\rangle \mapsto | \theta_i \rangle \left( \sqrt{1-\lambda_i}| 0 \rangle + 
\sqrt{\lambda_i}| 1 \rangle \right)$, and measure the ancilla in the 0,1
basis.}. By repeating this measurement on multiple copies of the same state,
we can estimate $\langle \tilde{O} \rangle$; for our purposes, it is enough
to have polynomially many copies of the state.

We now describe the verifier $V$. The witness $\tau$ consists of several
(i.e., polynomially many) blocks, where each block has $d$ qubits,
supposedly representing one copy of the state $\tilde{\sigma}$. On each
block, $V$ measures the observable $\sum_k | 1 \rangle \langle 1 |_k$, and
if the outcome does not equal $N$, $V$ outputs ``false.'' This projects each
block onto the space of $N$-particle states.

Next, $V$ performs measurements on each block, to estimate the expectation
values of $\tilde{\sigma}$, for a suitable set of observables. Then $V$
checks whether they match the expectation values of $\rho$. One problem
arises: the prover could try to cheat by entangling the different blocks of
qubits. One can show that this does not fool the verifier, using a Markov
argument, as was done in \cite{Aharonov}. This suffices to show that $N$%
-representability is in QMA, and hence finishes the proof.

What can be said about the complexity of the pure-state $N$%
-representability-problem, where one has to decide whether the reduced
density operators arise from a pure state of N fermions? In that case, the
verifier must be able to convince himself that the state he gets is pure.
This can be done when he gets two states $\rho$ and $\sigma$ that are
promised to be uncorrelated, i.e. that he gets the state $\rho\otimes\sigma$%
: then Arthur can calculate the expectation value of the observable $\Tr%
(\rho\sigma)$, and this can only be close to $1$ when $\rho$ is pure and $%
\sigma\simeq\rho$. Indeed, if $\Tr(\sigma^2) \leq 1-\varepsilon$, then for
all states $\tau$, $\Tr(\sigma\tau) \leq \sqrt{\Tr(\sigma^2)\Tr(\tau^2)}
\leq (1-\varepsilon)^{1/2} \leq 1-\varepsilon/2$. The problem however is
that Merlin can cheat, and hand out a correlated state to Arthur, such that
the above test becomes inconclusive. This is precisely the feature that
distinguishes the complexity class QMA(k) from QMA \cite{Matsumoto}: in
QMA(k), the verifier is promised to get a tensor product of different
states, and it has been conjectured that QMA(k) is strictly larger than QMA.
The above discussion shows that the pure-state $N$-representability problem
is contained in QMA(k), and it is hence plausible that it is harder than the 
$N$-representability problem. It would be very interesting to investigate
whether the problem is QMA(k)-complete.

%is also in QMA.  The idea is that the verifier can measure the purity of the
%state $\sigma$ using the ``swap test'':  assume we are given two copies of
%$\sigma$; adjoin an ancilla qubit prepared in the state $\ket{+} =
%(\ket{0}+\ket{1})/\sqrt{2}$; do a controlled-swap between the two copies of
%$\sigma$, controlled by the ancilla qubit; measure the ancilla qubit in the
%Hadamard basis $\ket{\pm} = (\ket{0}\pm\ket{1})/\sqrt{2}$.  Then the
%probability of observing $\ket{+}$ is $(1 + \Tr(\sigma^2))/2$, which
%measures the purity of $\sigma$.
%
%The main difficulty is that the swap test requires a product state $\sigma
%\otimes \sigma$, whereas a dishonest prover might give us an arbitrary state
%$\tau$.  To get around this, the verifier can do the following trick:  assume
%we are given an arbitrary state $\tau$, on $r+s$ blocks of qubits (where each
%block supposedly represents one copy of the state $\sigma$); randomly permute
%the blocks, then keep a subset of $r$ blocks, and discard the rest; let
%$\tau'$ denote the resulting state.  For our purposes, we will set $r =
%\poly(\n)$ and $s = \poly(r)$.
%
%By the de Finetti theorem for finite quantum states \cite{KR}, $\tau'$ is
%close to a mixture of $r$-fold tensor product states; in particular, there is
%a state $\tau'' = \int P(\sigma) \sigma^{\otimes r} d\sigma$ such that
%$\lVert \tau' - \tau'' \rVert_1 \leq O(r/\sqrt{s})$.  So the swap test will
%work properly on $\tau'$.  It is now straightforward to construct a QMA
%verifier for pure-state $\N$-representability; we omit the details.

%$\N$-representability also belong to the class QMA. This can be understood as follows. Consider a fermionic
%wavefunction (\ref{wave}), and the corresponding \N-qubit state
%\be|\psi_{qubit}\rangle=\sum_{{\tiny\begin{array}{l}i_1,i_2,... i_N=0\\ i_1+i_2+...=n \end{array}}}^1
%c_{i_1,i_2,...}|i_1\rangle|i_2\rangle...|i_N\rangle\label{wavequbit}\ee

%Every local expectation value on the original wavefunction $\langle\psi|a_i^\dagger a_j^\dagger a_k
%a_l|\psi\rangle$ can be mapped to an expectation value of a different operator on the state $|\psi_{qubit}\rangle$
%by doing a Jordan-Wigner transformation:
%\[ a_i\rightarrow \left(\otimes_{k<i}\hspace{.1cm}\sigma_k^z\right)|0_i\rangle\langle 1_i|\equiv A_i\]
%The operators $A_i$ are highly nonlocal, and we have

%\[\langle\psi|a_i^\dagger a_j^\dagger a_k a_l|\psi\rangle=\langle\psi_{qubit}|A_i^\dagger A_j^\dagger A_k
%A_l|\psi_{qubit}\rangle.\]

%Given a state $|\psi_{qubit}\rangle$, this expectation value can be evaluated efficiently, as it is just the
%expectation value of a tensor product of local observables. If Merlin would now give many copies of the state
%$|\psi_{qubit}\rangle$ to Arthur (note that he cannot gain anything by entangling the copies or giving different
%copies because of the argument given in cite{Aharonov,Yi-Kai}), Arthur could measure the corresponding observables
%and check consistency up to polynomial precision. This showes that $N$-representability is in QMA, finishing the
%proof.

\vskip 10pt

It is remarkable that checking consistency of 2-body reduced density
operators of fermionic states is so hard, while checking consistency of
1-body reduced density operators is simple \cite{Coleman}. This can be
easily understood from the previous discussion: the extreme points of the
convex set of 1-body density operators $\langle a_i^\dagger a_j\rangle$
correspond to ground states of Hamiltonians only containing bilinear terms
in $a_i^\dagger$ and $a_j$; such Hamiltonians can easily be diagonalized as
they represent systems of free fermions, and hence the consistency problem
can easily be solved. As shown in \cite{Coleman}, consistency can be decided
in that case based solely on the eigenvalues of the reduced density
operators. A number of related problems, where consistency only depends on
the eigenvalues, have been investigated recently~\cite{reduced}; most
remarkably, a characterisation has been obtained for the polytope of
1-particle marginals that are $N$-representable under the condition that the 
$N$-particle wavefunction must be pure \cite{Klyachko05}.

These results have to be contrasted with our problem of finding the $N$%
-representable 2-body density operators, where the eigenvalues alone are not
enough to decide consistency but also the eigenvectors are relevant.
Actually, let us consider the simpler problem of deciding $N$%
-representability of 2-fermion density operators where only the diagonal
elements $D_{ij}=\langle a_i^\dagger a_j^\dagger a_j a_i\rangle$ are
specified. If we consider the case $d=2N$ and the mapping discussed above,
one easily finds that the extreme points of this set would be obtained by
ground states of local spin Hamiltonians which only contain commuting $%
\sigma^z$ operators. These correspond to spin-glasses, and so the problem of
deciding $N$-representability of $\{D_{ij}\}$ is NP-hard \cite{Barahona}. It
was indeed pointed out a long time ago that $N$-representability restricted
to the diagonal elements is equivalent to a combinatorial problem \cite{Kuhn}
that was later shown to be equivalent to the NP-hard problem of deciding
membership in the boolean quadric polytope \cite{Deza}.

Let's finally discuss the relevance of the above results in the context of
quantum chemistry. We have shown that it is a hopeless task to determine the
ground states of all local fermionic Hamiltonians, and in particular, that
an approach by means of the $N$-representability problem is intractable,
even on a quantum computer. But it is possible that other physical systems,
e.g.~systems with different particle statistics, additional symmetry or a
limited number of modes, might allow for an efficient characterisation of
the two-particle reduced density matrices, and hence an efficient
calculation of the ground states of local Hamiltonians. This seems to be the
case in e.g. one-dimensional translational invariant spin systems, where the
density matrix renormalization group \cite{DMRG} allows for a systematic
approximation of the allowed convex set of reduced density operators from
within \cite{VC06}. A very different numerical approach has been developed
in the context of quantum chemistry and it known as the the contracted Schr{%
\"o}dinger equation~\cite{book}. In that method, one approximates the convex
set from the outside, and the $N$-representability problem is a crucial
ingredient of the algorithm. The foregoing discussion shows that this
approach has to break down in the most general case, and it would be very
interesting to investigate the conditions under which those approximations
are justified.

%The results in this paper show that it is a hopeless task to
%determine ground state of all local fermionic Hamiltonians. But it
%could well be that this problem becomes tractable when there is more
%symmetry present.

%Another interesting problem is to characterize the
%convex set of reduced density operators of a translational invariant
%state (with bounded local dimension). There are very good
%indications that the complexity of this problem is polynomial, but a
%proof is still missing \cite{VCfaith}. Furthermore, if the system is
%permutation invariant, an efficient characterisation of the
%two-particle reduced density matrix is possible, because
%$\N$-particle permutation symmetric states can be characterised with
%a polynomial number of parameters using Schur-Weyl duality.
%%This follows from Schur-Weyl duality which allows us to rewrite
%%$$ \sigma=\sum_{\lambda} \omega_\lambda \sigma_{U_\lambda} \otimes \frac{\openone_{V_\lambda}}{\textrm{dim} V_\lambda}.$$
%%The sum extends over all partitions $\lambda$ of $\N$ into $d$ parts
%%and $U_\lambda$ ($V_\lambda$) is the space of the representation of
%%the unitary group $U(d)$ (the symmetric group $S_n$) with highest
%%weight (Young frame) $\lambda$. The number of free parameters is a
%%polynomial in $\N$ since $\dim U_\lambda$ as well as the number of
%%partitions are both bounded by polynomials in $\N$.
%However, this argument breaks down for growing $d$.
%%And indeed, as we saw before $d$ cannot grow arbitrarily fast,
%%else we could solve $\N$-representability.
%
%However, with the help of quantum de Finetti
%theorems we can obtain results for slowly growing $d$. Such theorems
%quantify how well $k$-particle reduced density matrices arising from
%permutation-invariant states can be approximated by convex
%combinations of tensor product states~\cite{Hudson, CKMR06}. The set
%of tensor product density operators can be efficiently characterised
%and (for fixed $k$) the approximation is still an inverse polynomial
%in $\poly(\N)$ if $d=o(\N^{\frac{1}{2}-\epsilon})$ for $\epsilon>0$
%and in the case of bosons even if
%$d=o(\N^{1-\epsilon})$~\cite{CKMR06}.
%(Note that the remarkable success of (dynamical) mean field theories
%relies on the fact that tensor product density operators are often
%very good approximations of the real ones, despite the fact that the
%underlying Hamiltonian has no permutation symmetry.)
%
%%The intuitive reason for the efficiency can in this case be traced
%%to so-called de Finetti theorems.
%
%%$\N$-representability can be solved when the reduced density
%%operators are so-called "exchangable". This happens when the
%%underlying Hamiltonian is permutation invariant with respect to the
%%modes and when additionally the number of modes tends to infinity.
%%In that case, one can use the quantum de-Finetti theorem
%%\cite{Hudson} (also valid in the case of fermionic statistics!) and
%%infer that the corresponding extreme points of the 2-particle
%%density operators are tensor products of the extreme 1-particle
%%density operators.

In conclusion, we investigated the problem of $N$-representability, and
characterized its computational complexity by showing that it is
QMA-complete. Obviously, the theory of quantum computing was a prerequisite
to pinpoint the computational complexity of this classic problem. .

\vskip 10pt

\noindent \textit{Acknowledgements:} Y.K.L. and M.C. thank the Institute for
Quantum Information for its hospitality. Y.K.L. is supported by an ARO/DTO
Quantum Computing Graduate Research Fellowship. M.C. acknowledges an EPSRC
Postdoctoral and a Nevile Research Fellowship which he holds at Magdalene
College Cambridge, and is supported by the EU under the FP6-FET Integrated
Project SCALA, CT-015714. F.V. is supported by the Gordon and Betty Moore
Foundation through Caltech's Center for the Physics of Information, and by
the National Science Foundation under Grant No. PHY-0456720.

\begin{thebibliography}{99}
\bibitem{} 
%\bibitem{CKMR06} M. Christandl, R. K{\"o}nig, G. Mitchison and R. Renner,
%quant-ph/0602130 (2006).

\bibitem{Coulson} C. A. Coulson, Rev. Mod. Phys. \textbf{32}, 175 (1960); R.
H. Trethold, Phys. Rev. \textbf{105}, 1421 (1957).

\bibitem{Coleman} A. J. Coleman, Rev. Mod. Phys. \textbf{35}, 668 (1963).

\bibitem{book1} A.J. Coleman and V.I. Yukalov, \textit{Reduced Density
Matrices: Coulson's Challenge}, Springer (2000).

\bibitem{book} J. Cioslowski, \textit{Many-Electron Densities and Reduced
Density Matrices}, Kluwer Academic (2000).

\bibitem{Mazziotti} D. A. Mazziotti, Acc. Chem. Res. \textbf{39}, 207 (2006).

\bibitem{Kitaev} A. Yu Kitaev, A. H. Shen and M. N. Vyalyi, \textit{%
Classical and Quantum Computation}, AMS (2002).

\bibitem{Kempe1} J. Kempe and O. Regev, Quant. Info. and Comput. \textbf{3},
258 (2003).

\bibitem{Kempe2} J. Kempe, A. Kitaev and O. Regev, SIAM J. Comput. \textbf{35%
}, 1070 (2006); quant-ph/0406180.

\bibitem{Terhal} R. Oliveira and B. Terhal, quant-ph/0504050.

\bibitem{Yi-Kai} Y.-K. Liu, Proc. RANDOM 2006, 438; quant-ph/0604166.

\bibitem{GLS} M. Gr\"otschel, L. Lov\'asz and A. Schrijver, \textit{%
Geometric Algorithms and Combinatorial Optimization}, Springer-Verlag (1988).

\bibitem{Barahona} F. Barahona, J. Phys. A: Math. Gen. \textbf{15}, 3241
(1982).

\bibitem{VC} F. Verstraete and J.I. Cirac, J. Stat. Mech. P09012 (2005).

\bibitem{BV} D. Bertsimas and S. Vempala, J. Assoc. Comput. Mach. \textbf{51}%
, 540-556 (2004).

\bibitem{Aharonov} D. Aharonov and O. Regev, Proc. FOCS 2003, 210;
quant-ph/0307220. 
%\bibitem{KR} R. K\"onig and R. Renner, J. Math. Phys. 46, 122108 (2005); quant-ph/0410229.

\bibitem{Matsumoto} H. Kobayashi, K. Matsumoto and Tomoyuki Yamakami,
quant-ph/0110006.

\bibitem{reduced} A. Higuchi, A. Sudbery, J. Szulc, Phys. Rev. Lett. \textbf{%
90}, 107902 (2003); S. Bravyi, Quantum Inf. and Comp. \textbf{4}, 12 (2004);
M. Christandl and G. Mitchison, Commun. Math. Phys. \textbf{261}, 789
(2006); A. Klyachko, quant-ph/0409113; S. Daftuar and P. Hayden, Ann. Phys. 
\textbf{315}, 80 (2005).

\bibitem{Klyachko05} A. Klyachko, J. of Physics: Conference Series \textbf{36%
} (1) pp.72-86 (2006).

\bibitem{Kuhn} M. L. Yoseloff and H. W. Kuhn, Journ. of Math. Phys. \textbf{%
10}, 703 (1969).

\bibitem{Deza} M. Deza and M. Laurent, Journ. of Comp. and Appl. Math. 
\textbf{55}, 217 (1994).

\bibitem{DMRG} U. Schollwoeck, Rev. Mod. Phys. \textbf{77}, 259 (2005).

\bibitem{VC06} F. Verstraete and J.I. Cirac, Phys. Rev. B \textbf{73},
094423 (2006). 
%\bibitem{Hudson}  R. L. Hudson and G. R. Moody, Wahrscheinlichkeitstheorie
%Gebiete {\bf 33}, 343 (1976). C. M. Caves, C. A. Fuchs and R.
%Schack, J. Math. Phys. {\bf 43}, 4537 (2002). R. K{\"o}nig and R
%Renner, J. Math. Phys. {\bf 46}, 122108 (2005). R. Renner, PhD
%thesis, ETH Z{\"u}rich, quant-ph/0512258 (2005).
\end{thebibliography}

\end{document}
